%% =====================================================================
%%  Bilan 1 — Premiers contacts (unités 1 à 4)
%%  Pistes : B1.01 (écoute notée) · B1.02 (expression orale)
%% =====================================================================
\bilan{1}

\noindent{\fontsize{11}{14.6}\selectfont Vous avez terminé la première partie : bravo ! Ce bilan
vous permet de vérifier vos acquis. Révisez d'abord l'essentiel des unités 1 à 4, puis faites
les huit exercices sans regarder les unités : ils sont notés sur 20. Les corrigés sont en fin de
livre, page~\pageref{cor-b1-1}.\par}

%% ---------------------------------------------------------------------
\rubrique{grammaire}{Révisez}

\begin{paire}
\begin{regle}[Le verbe \en{be} · unités 1 et 2]
\renewcommand{\arraystretch}{1.25}
\begin{tabular}{@{}>{\enOn}l<{\enOff}@{\hspace{3mm}}>{\enOn}l<{\enOff}@{\hspace{3mm}}>{\enOn}l<{\enOff}@{}}
\otTH{Affirmation} & \otTH{Négation} & \otTH{Question}\\
I'm      & I'm not      & Am I…?\\
he's     & he isn't     & Is he…?\\
she's    & she isn't    & Is she…?\\
it's     & it isn't     & Is it…?\\
we're    & we aren't    & Are we…?\\
you're   & you aren't   & Are you…?\\
they're  & they aren't  & Are they…?\\
\end{tabular}\par\vspace{1.4mm}
Réponses courtes : \en{Yes, I am. / No, I'm not.}
\end{regle}
\begin{regle}[Les questions · unités 2 et 3]
\exen{What's your name?}{nom, numéro}
\exen{Where are you from?}{lieu}
\exen{Who's that?}{personne}
\exen{How are you?}{état}
\exen{How old are you?}{âge}
\exen{How do you spell it?}{orthographe}
\par\vspace{1mm}L'âge se dit avec \en{be} : \en{I'm twenty.}
\end{regle}
\end{paire}

\begin{paire}
\begin{regle}[Les possessifs · unité 4]
\en{my · your · his · her · its · our · your · their}\par\vspace{1mm}
\en{his} (à lui) ou \en{her} (à elle) : selon le \textbf{possesseur}.\par\vspace{1mm}
\exen{Léa and her brother}{Léa et son frère}
\end{regle}
\begin{regle}[\en{a} / \en{an} et le génitif · unités 3 et 4]
\en{a phone}, \en{an email}, \en{an hour} : selon le \textbf{son} qui suit.\par\vspace{1mm}
\exen{Léa's brother}{le frère de Léa}
\exen{my parents' house}{la maison de mes parents}
\end{regle}
\end{paire}

%% ---------------------------------------------------------------------
\rubrique{pratique}{Testez-vous}

\exercice[ecrire,niveau=1,points=3]{Le verbe \en{be}}{Complétez avec la forme qui convient :
affirmative ou négative, selon le sens. \textit{(0,5 point par réponse)}}
\begin{questions}
\item Léa \trou[16mm]{is} French. She \trou[16mm]{isn't} American.
\item Jack and Emma \trou[16mm]{are} Léa's colleagues.
\item I \trou[16mm]{am} a beginner, but my English \trou[16mm]{is} OK!
\item Mrs.\ Lopez \trou[16mm]{isn't} Mexican: she's American.
\end{questions}

\exercice[ecrire,niveau=2,points=3]{Questions et réponses}{Écrivez la question (1 à 4) ou la réponse
courte (5 et 6). \textit{(0,5 point par réponse)}}
\begin{questions}
\item \trou[60mm]{Where are you from?} — I'm from Lyon.
\item \trou[60mm]{How old are you?} — I'm twenty-seven.
\item \trou[60mm]{What's your name?} — My name is Sam.
\item \trou[60mm]{Who's that?} — That's my sister, Maya.
\item Is Emma British? — \trou[50mm]{Yes, she is.}
\item Are Jack and Emma French? — \trou[50mm]{No, they aren't.}
\end{questions}

\exercice[ecrire,niveau=1,points=2]{\en{his}, \en{her} ou \en{their} ?}{Complétez.
\textit{(0,5 point par réponse)}}
\begin{questions}
\item Sam is here with \trou[18mm]{his} sister, Maya.
\item Mrs.\ Lopez is with \trou[18mm]{her} husband, Miguel.
\item Sofia and David are in Chicago with \trou[18mm]{their} sons.
\item Léa is on the phone with \trou[18mm]{her} brother.
\end{questions}

\exercice[ecrire,niveau=2,points=2]{Le génitif}{Traduisez avec \en{'s} ou \en{'}.
\textit{(0,5 point par réponse)}}
\begin{questions}[2]
\item le chat de Sam → \trou[40mm]{Sam's cat}
\item la fille de Mrs.\ Lopez → \trou[40mm]{Mrs. Lopez's daughter}
\item le frère d'Emma → \trou[40mm]{Emma's brother}
\item les noms des garçons → \trou[40mm]{the boys' names}
\end{questions}

\exercice[ecrire,niveau=1,points=2]{Les nombres}{Écrivez les nombres en lettres.
\textit{(0,5 point par réponse)}}
\begin{questions}[4]
\item 14 → \trou[24mm]{fourteen}
\item 40 → \trou[24mm]{forty}
\item 19 → \trou[24mm]{nineteen}
\item 72 → \trou[24mm]{seventy-two}
\end{questions}

\exercice[ecouter,niveau=2,points=3,piste=B1.01]{À la fête de Sam}{Samedi soir, Léa et Emma arrivent
à la fête de Sam. Écoutez, puis indiquez si chaque phrase est vraie (V) ou fausse (F).
\textit{(0,5 point par réponse)}}
\begin{questions}[2]
\item Emma is Sam's neighbor. \hfill\trou[9mm]{F}
\item Emma is British. \hfill\trou[9mm]{V}
\item Emma is from Manchester. \hfill\trou[9mm]{F}
\item Maya is Sam's sister. \hfill\trou[9mm]{V}
\item Maya is twenty-two. \hfill\trou[9mm]{F}
\item It's Sam's party. \hfill\trou[9mm]{V}
\end{questions}

\exercice[lire,niveau=2,points=3]{Le profil de Sam}{Lisez le profil de Sam, puis indiquez si chaque
phrase est vraie (V) ou fausse (F). \textit{(0,5 point par réponse)}}
\begin{anglais}
\begin{tcolorbox}[otBase,colback=white,colframe=otRule,boxrule=0.5pt,fontupper=\fontsize{10.8}{14.4}\selectfont]
\begin{minipage}[t]{20mm}\vspace{0pt}\otAvatar[otTealTint]{sam}{18mm}\end{minipage}\hfill
\begin{minipage}[t]{\dimexpr\linewidth-24mm\relax}\vspace{0pt}
Hi! I'm Sam Brooks. I'm twenty-four and I'm from Brooklyn, New York. I'm a music student and a
barista. My parents are in Atlanta, but my sister, Maya, is in New York with me. She's twenty.
Our cat's name is Mango. My phone number? Sorry, it's a secret!
\par\vspace{1mm}{\fontsize{8.6}{10}\selectfont\color{otMuted}\foreignlanguage{french}{\textit{a music student} : un étudiant en musique · \textit{a barista} : un serveur de café · \textit{a secret} : un secret}}
\end{minipage}
\end{tcolorbox}
\begin{questions}[2]
\item Sam is twenty-four. \hfill\trou[9mm]{V}
\item Sam is from Atlanta. \hfill\trou[9mm]{F}
\item Sam's parents are in Atlanta. \hfill\trou[9mm]{V}
\item Maya is Sam's sister. \hfill\trou[9mm]{V}
\item Mango is Sam's dog. \hfill\trou[9mm]{F}
\item Sam's phone number is in his profile. \hfill\trou[9mm]{F}
\end{questions}
\end{anglais}

\exercice[ecrire,niveau=3,points=2]{Votre présentation}{Vous écrivez à un nouveau collègue. En cinq
phrases, présentez-vous : nom, âge, origine, langues parlées, famille.
\textit{(0,5 point par information correcte, 2 points au maximum)}}
\lignes{5}

%% ---------------------------------------------------------------------
\rubrique{ecouter}{Parlez}

\exercice[ecouter,parler,niveau=3,piste=B1.02]{Rencontre à la fête}{À la fête de Sam, Maya vient vous
parler. Répondez-lui \textbf{à voix haute} pendant chaque pause. Cet exercice n'est pas noté :
comparez vos réponses avec les modèles du corrigé.}
\begin{center}\otTableau
\begin{tabular}{@{}>{\otSemi\fontsize{8.6}{10}\selectfont\color{otNavy}}l@{\hspace{3mm}}>{\enOn}p{70mm}<{\enOff}@{\hspace{6mm}}p{78mm}@{}}
 & \otTH{Maya dit\ldots} & \otTH{Vous\ldots}\\\midrule
1. & Hi! I'm Maya, Sam's sister. What's your name? & saluez-la et donnez votre prénom.\\
2. & Nice to meet you! Where are you from? & réagissez, puis dites d'où vous venez.\\
3. & Oh, cool! Are you a friend of Léa? & répondez par une réponse courte.\\
4. & And is your family in New York? & répondez et dites où est votre famille.\\
5. & OK! Welcome to the party! & remerciez-la.\\\bottomrule
\end{tabular}\end{center}

\otScore{20}
